\subsection{Notation and organization}
\begin{table}[h]\centering\caption{Notation used throughout.}
\begin{tabular}{ll}\hline
Symbol & Meaning \\ \hline
$s$, $\ell$ & a peptide sequence and its length \\
$L$ & encoding cap (150; pilot 60; full Feb2020 audit 200) \\
$K_k(s)$ & set of contiguous $k$-mers of $s$ \\
$E(s)$ & encoded sequence, $\mathbb{R}^{L\times(20+d)}$ \\
$\tilde A$ & symmetric-normalized shared adjacency \\
$\hat p$, $f(s)$ & sigmoid classifier output and ensemble score \\
$\omega$ & positive-class weight $n_-/n_+$ \\
$\pi_T(s)$, $T$ & fixed-temperature Boltzmann measure and temperature \\
$J(\cdot,\cdot)$ & Jaccard index \\
$B$ & resampling count (bootstrap or permutation) \\
\hline\end{tabular}\end{table}

Section 2 situates the work; Section 3 defines the data, clustering, and
splits; Section 4 defines the models, objective, designer, and statistical
procedures, with proofs; Section 5 fixes the evaluation metrics; Later sections report the benchmark studies, candidate screen, signal
ablation, threats, reproducibility audit and software artifact. Appendices
give hyperparameters, provenance and design/screen listings.
